% 連続体力学 連続体力学-1 問題1（001-temperature） 解いた日 2026-10-17
%% kotae
\toi{1}{◎}{}
\[ \pdv{T}{t}=-\frac{a}{\tau}x_{1}^{2}\ee^{-t/\tau} \]
\toi{2}{◎}{}
\begin{align*}
(\vect{u}\cdot\vect{\nabla})T&=kx_{1}\pdv{}{x_{1}}T-kx_{2}\pdv{}{x_{2}}T+0\\
&=kx_{1}(2ax_{1}\ee^{-t/\tau})-kx_{2}\cdot0\\
&=2akx_{1}^{2}\ee^{-t/\tau}
\end{align*}
\toi{3}{◎}{}
\begin{align*}
\frac{\mathrm{D}T}{\mathrm{D}t}&=\pdv{T}{t}+(\vect{u}\cdot\vect{\nabla})T\\
&=-\frac{a}{\tau}x_{1}^{2}\ee^{-t/\tau}+2akx_{1}^{2}\ee^{-t/\tau}\\
&=\Bigl(-\frac{1}{\tau}+2k\Bigr)ax_{1}^{2}\ee^{-t/\tau}
\end{align*}
\toi{4}{○}{記述}
\[ -\frac{1}{\tau}+2k=0,\quad \therefore\ \tau=\frac{1}{2k} \]
温度が低い粒子が流れこんでくるから。\nuke{粒子の動き}
\naze{問いは「粒子の動きに注目して」説明することを求めている。粒子は $u_{1}=kx_{1}$ により $x_{1}=0$ から遠ざかる向きに流され（$x_{1}(t)=X_{1}\ee^{kt}$），各粒子の温度は $T_{0}+aX_{1}^{2}$ で一定である。時刻 $t$ に固定点 $x_{1}$ にいる粒子は $\abs{X_{1}}=\abs{x_{1}}\ee^{-kt}$ から来たので，後から来る粒子ほど $x_{1}=0$ に近い（温度の低い）位置の出身である，まで書く。}
%% kizuki
\item 「粒子の動きに注目して説明せよ」では，粒子の軌道 $x_{1}(t)=X_{1}\ee^{kt}$ を求め，粒子の温度を $X_{1}$ で表し（$T_{0}+aX_{1}^{2}$ で一定），固定点に来る粒子の出発点 $\abs{X_{1}}=\abs{x_{1}}\ee^{-kt}$ が時間とともに $x_{1}=0$ に近づく，という順に式で示す。「低い粒子が流れ込む」は，どこから来た粒子かを書いて初めて理由になる。
